\documentclass[11pt]{article}

\usepackage{fullpage}
\usepackage{newtxtext}
\usepackage{microtype}
\usepackage[T1]{fontenc}

\usepackage{hyperref}
\hypersetup{hypertexnames=false,breaklinks=true,colorlinks=true,linkcolor=blue,filecolor=blue,citecolor=blue,urlcolor=blue}

\usepackage{mathtools}
% show equation numbers only for referenced equations
\mathtoolsset{showonlyrefs}

\usepackage{amsfonts,amssymb}
\usepackage[amsmath,amsthm,thmmarks,hyperref]{ntheorem}
\usepackage{thmtools}

\usepackage{enumitem}

\usepackage{zref-clever}
\zcsetup{
        cap=true, % capitalize types
        nameinlink,
}
% for migration from cleveref
\newcommand{\cref}[1]{\zcref{#1}}
\newcommand{\Cref}[1]{\zcref[S]{#1}}

%%% BLACKBOARD SYMBOLS

\newcommand{\C}{\ensuremath{\mathbb{C}}}
\newcommand{\D}{\ensuremath{\mathbb{D}}}
\newcommand{\F}{\ensuremath{\mathbb{F}}}
\newcommand{\G}{\ensuremath{\mathbb{G}}}
\newcommand{\J}{\ensuremath{\mathbb{J}}}
\newcommand{\N}{\ensuremath{\mathbb{N}}}
\newcommand{\Q}{\ensuremath{\mathbb{Q}}}
\newcommand{\R}{\ensuremath{\mathbb{R}}}
\newcommand{\T}{\ensuremath{\mathbb{T}}}
\newcommand{\Z}{\ensuremath{\mathbb{Z}}}
\newcommand{\QR}{\ensuremath{\mathbb{QR}}}

\newcommand{\Zt}{\ensuremath{\Z_t}}
\newcommand{\Zp}{\ensuremath{\Z_p}}
\newcommand{\Zq}{\ensuremath{\Z_q}}
\newcommand{\ZN}{\ensuremath{\Z_N}}
\newcommand{\Zps}{\ensuremath{\Z_p^*}}
\newcommand{\ZNs}{\ensuremath{\Z_N^*}}
\newcommand{\JN}{\ensuremath{\J_N}}
\newcommand{\QRN}{\ensuremath{\QR_{N}}}
\newcommand{\QRp}{\ensuremath{\QR_{p}}}

%%% THEOREM COMMANDS

\declaretheorem[numberwithin=section]{theorem}

\declaretheorem[sibling=theorem,style=plain]
        {lemma,corollary,proposition,claim,fact}

\declaretheorem[sibling=theorem,style=definition]
        {definition,conjecture,construction,assumption,protocol}

\declaretheorem[sibling=theorem,style=remark]
        {remark,example,note}

% teach zref-clever these names

\zcRefTypeSetup{claim}{Name-sg=Claim,Name-pl=Claims,name-sg=claim,name-pl=claims}
\zcRefTypeSetup{fact}{Name-sg=Fact,Name-pl=Facts,name-sg=fact,name-pl=facts}
\zcRefTypeSetup{maintheorem}{Name-sg=Main Theorem,Name-pl=Main
Theorems,name-sg=main theorem,name-pl=main theorems}
\zcRefTypeSetup{conjecture}{Name-sg=Conjecture,Name-pl=Conjectures,name-sg=conjecture,name-pl=conjectures}
\zcRefTypeSetup{construction}{Name-sg=Construction,Name-pl=Constructions,name-sg=construction,name-pl=constructions}
\zcRefTypeSetup{assumption}{Name-sg=Assumption,Name-pl=Assumptions,name-sg=assumption,name-pl=assumptions}
\zcRefTypeSetup{protocol}{Name-sg=Protocol,Name-pl=Protocols,name-sg=protocol,name-pl=protocols}
\zcRefTypeSetup{question}{Name-sg=Question,Name-pl=Questions,name-sg=question,name-pl=questions}

% equation numbering style
\numberwithin{equation}{section}

%%% GENERAL COMPUTING

\newcommand{\bit}{\ensuremath{\set{0,1}}}
\newcommand{\pmone}{\ensuremath{\set{-1,1}}}

% asymptotics
\DeclareMathOperator{\poly}{poly}
\DeclareMathOperator{\polylog}{polylog}
\DeclareMathOperator{\negl}{negl}
\newcommand{\Otil}{\ensuremath{\tilde{O}}}

% probability/distribution stuff
\DeclareMathOperator*{\E}{E}
\DeclareMathOperator*{\Var}{Var}

% sets in calligraphic type
\newcommand{\calD}{\ensuremath{\mathcal{D}}}
\newcommand{\calF}{\ensuremath{\mathcal{F}}}
\newcommand{\calG}{\ensuremath{\mathcal{G}}}
\newcommand{\calH}{\ensuremath{\mathcal{H}}}
\newcommand{\calX}{\ensuremath{\mathcal{X}}}
\newcommand{\calY}{\ensuremath{\mathcal{Y}}}

% types of indistinguishability
\newcommand{\compind}{\ensuremath{\stackrel{c}{\approx}}}
\newcommand{\statind}{\ensuremath{\stackrel{s}{\approx}}}
\newcommand{\perfind}{\ensuremath{\equiv}}

% font for general-purpose algorithms
\newcommand{\algo}[1]{\ensuremath{\mathsf{#1}}}
% font for general-purpose computational problems
\newcommand{\problem}[1]{\ensuremath{\mathsf{#1}}}
% font for complexity classes
\newcommand{\class}[1]{\ensuremath{\mathsf{#1}}}

% complexity classes and languages
\renewcommand{\P}{\class{P}}
\newcommand{\BPP}{\class{BPP}}
\newcommand{\NP}{\class{NP}}
\newcommand{\coNP}{\class{coNP}}
\newcommand{\AM}{\class{AM}}
\newcommand{\coAM}{\class{coAM}}
\newcommand{\IP}{\class{IP}}

%%% "LEFT-RIGHT" PAIRS OF SYMBOLS


% inner product
\DeclarePairedDelimiter\inner{\langle}{\rangle}
% absolute value
\DeclarePairedDelimiter\abs{\lvert}{\rvert}
% a set
\DeclarePairedDelimiter\set{\{}{\}}
% parens
\DeclarePairedDelimiter\parens{(}{)}
% tuple, alias for parens
\DeclarePairedDelimiter\tuple{(}{)}
% square brackets
\DeclarePairedDelimiter\bracks{[}{]}
% rounding off
\DeclarePairedDelimiter\round{\lfloor}{\rceil}
% floor function
\DeclarePairedDelimiter\floor{\lfloor}{\rfloor}
% ceiling function
\DeclarePairedDelimiter\ceil{\lceil}{\rceil}
% length of some vector, element
\DeclarePairedDelimiter\length{\lVert}{\rVert}
% "lifting" of a residue class
\DeclarePairedDelimiter\lift{\llbracket}{\rrbracket}
\DeclarePairedDelimiter\len{\lvert}{\rvert}

%%% CRYPTO-RELATED NOTATION

% KEYS AND RELATED

\newcommand{\key}[1]{\ensuremath{#1}}

\newcommand{\pk}{\key{pk}}
\newcommand{\vk}{\key{vk}}
\newcommand{\sk}{\key{sk}}
\newcommand{\mpk}{\key{mpk}}
\newcommand{\msk}{\key{msk}}
\newcommand{\fk}{\key{fk}}
\newcommand{\id}{id}
\newcommand{\keyspace}{\ensuremath{\mathcal{K}}}
\newcommand{\msgspace}{\ensuremath{\mathcal{M}}}
\newcommand{\ctspace}{\ensuremath{\mathcal{C}}}
\newcommand{\tagspace}{\ensuremath{\mathcal{T}}}
\newcommand{\idspace}{\ensuremath{\mathcal{ID}}}

\newcommand{\concat}{\ensuremath{\|}}

% GAMES

% advantage
\newcommand{\advan}{\ensuremath{\mathbf{Adv}}}

% different attack models
\newcommand{\attack}[1]{\ensuremath{\text{#1}}}

\newcommand{\atk}{\attack{atk}} % dummy attack
\newcommand{\indcpa}{\attack{ind-cpa}}
\newcommand{\indcca}{\attack{ind-cca}}
\newcommand{\anocpa}{\attack{ano-cpa}} % anonymous
\newcommand{\anocca}{\attack{ano-cca}}
\newcommand{\euacma}{\attack{eu-acma}} % forgery: adaptive chosen-message
\newcommand{\euscma}{\attack{eu-scma}} % forgery: static chosen-message
\newcommand{\suacma}{\attack{su-acma}} % strongly unforgeable

% ADVERSARIES
\newcommand{\attacker}[1]{\ensuremath{\mathcal{#1}}}

\newcommand{\Adv}{\attacker{A}}
\newcommand{\AdvA}{\attacker{A}}
\newcommand{\AdvB}{\attacker{B}}
\newcommand{\Dist}{\attacker{D}}
\newcommand{\Sim}{\attacker{S}}
\newcommand{\Ora}{\attacker{O}}
\newcommand{\Inv}{\attacker{I}}
\newcommand{\For}{\attacker{F}}

% CRYPTO SCHEMES

\newcommand{\scheme}[1]{\ensuremath{\text{#1}}}

% pseudorandom stuff
\newcommand{\prg}{\algo{PRG}}
\newcommand{\prf}{\algo{PRF}}
\newcommand{\prp}{\algo{PRP}}

% symmetric-key cryptosystem
\newcommand{\skc}{\scheme{SKC}}
\newcommand{\skcgen}{\algo{Gen}}
\newcommand{\skcenc}{\algo{Enc}}
\newcommand{\skcdec}{\algo{Dec}}

% public-key cryptosystem
\newcommand{\pkc}{\scheme{PKC}}
\newcommand{\pkcgen}{\algo{Gen}}
\newcommand{\pkcenc}{\algo{Enc}} % can also use \kemenc and \kemdec
\newcommand{\pkcdec}{\algo{Dec}}

% digital signatures
\newcommand{\sig}{\scheme{SIG}}
\newcommand{\siggen}{\algo{Gen}}
\newcommand{\sigsign}{\algo{Sign}}
\newcommand{\sigver}{\algo{Ver}}

% message authentication code
\newcommand{\mac}{\scheme{MAC}}
\newcommand{\macgen}{\algo{Gen}}
\newcommand{\mactag}{\algo{Tag}}
\newcommand{\macver}{\algo{Ver}}

% key-encapsulation mechanism
\newcommand{\kem}{\scheme{KEM}}
\newcommand{\kemgen}{\algo{Gen}}
\newcommand{\kemenc}{\algo{Encaps}}
\newcommand{\kemdec}{\algo{Decaps}}

% identity-based encryption
\newcommand{\ibe}{\scheme{IBE}}
\newcommand{\ibesetup}{\algo{Setup}}
\newcommand{\ibeext}{\algo{Ext}}
\newcommand{\ibeenc}{\algo{Enc}}
\newcommand{\ibedec}{\algo{Dec}}

% hierarchical IBE (as key encapsulation)
\newcommand{\hibe}{\scheme{HIBE}}
\newcommand{\hibesetup}{\algo{Setup}}
\newcommand{\hibeext}{\algo{Extract}}
\newcommand{\hibeenc}{\algo{Encaps}}
\newcommand{\hibedec}{\algo{Decaps}}

% binary tree encryption (as key encapsulation)
\newcommand{\bte}{\scheme{BTE}}
\newcommand{\btesetup}{\algo{Setup}}
\newcommand{\bteext}{\algo{Extract}}
\newcommand{\bteenc}{\algo{Encaps}}
\newcommand{\btedec}{\algo{Decaps}}

% trapdoor functions
\newcommand{\tdf}{\scheme{TDF}}
\newcommand{\tdfgen}{\algo{Gen}}
\newcommand{\tdfeval}{\algo{Eval}}
\newcommand{\tdfinv}{\algo{Invert}}
\newcommand{\tdfver}{\algo{Ver}}

%%% PROTOCOLS

\newcommand{\out}{\text{out}}
\newcommand{\view}{\text{view}}

%%% COMMANDS FOR LECTURES/HOMEWORKS

\RequirePackage{fancyhdr}

\newcommand{\lecheader}{%
  \chead{\large \textbf{Lecture \lecturenum\\\lecturetopic}}

  \lhead{\small
    \textbf{\href{https://github.com/cpeikert/TheoryOfCryptography}{Theory
        of Cryptography}}\\}

  \rhead{\small \textbf{Instructor:
      \href{http://www.eecs.umich.edu/~cpeikert/}{Chris Peikert}\\Scribe:
      \scribename}}

  \setlength{\headheight}{20pt}
  \setlength{\headsep}{16pt}
}


%%% xsim settings

\RequirePackage{xsim}
\RequirePackage{needspace}

\DeclareExerciseEnvironmentTemplate{QRef}{%
    \par\addvspace{\baselineskip}
    \Needspace*{2\baselineskip}
    \noindent
    \hyperref[ans:\GetExerciseProperty{ID}]{\textbf{\XSIMmixedcase{\GetExerciseName}~\GetExerciseProperty{counter}.}}
    \label{ques:\GetExerciseProperty{ID}}
}{\par\addvspace{\baselineskip}}

\DeclareExerciseEnvironmentTemplate{ARef}{%
        \Needspace*{2\baselineskip}
        \noindent
        \hyperref[ques:\GetExerciseProperty{ID}]{\textbf{\XSIMmixedcase{\GetExerciseParameter{exercise-name}}~\GetExerciseProperty{counter}.}}
        \label{ans:\GetExerciseProperty{ID}}
        \GetExerciseBody{exercise}
        \par \medskip \noindent
        \textbf{\XSIMmixedcase{\GetExerciseName}.}
}{\par\addvspace{\baselineskip}}

\DeclareExerciseHeadingTemplate{simple}{%
    \section*{\XSIMmixedcase{\GetExerciseParameter{solution-name}s}}
}

\DeclareExerciseType{question}{%
    exercise-env = question,
    solution-env = answer,
    exercise-name = Question,
    solution-name = Answer,
    exercise-template = QRef,
    solution-template = ARef,
}

\xsimsetup{%
  path = {questions},
  file-extension = {aux},
  print-solutions/headings-template=simple,
  exercise/template=runin
}

\loadxsimstyle{layouts}

\AtEndDocument{\pagebreak\printsolutions}


\usepackage{algorithm,algorithmic}

% VARIABLES

\newcommand{\lecturenum}{5}
\newcommand{\lecturetopic}{Indistinguishability, Pseudorandomness}
\newcommand{\scribename}{Abhinav Shantanam}

% END OF VARIABLES

\lecheader

\pagestyle{plain}               % default: no special header

\begin{document}

\thispagestyle{fancy} % first page should have special header

% LECTURE MATERIAL STARTS HERE

\section{Indistinguishability}
\label{sec:indistinguishability}

We'll now start a major unit on \emph{indistinguishability} and
\emph{pseudorandomness}.  These concepts are a cornerstone of modern
cryptography, underlying several foundational applications such as
pseudorandom generators, secure encryption, ``commitment'' schemes,
and much more.

For example, our most immediate application of indistinguishability
will be to construct cryptographically strong \emph{pseudorandom bit
  generators}.  These are algorithms that produce many
``random-looking'' bits, while using very little ``true'' randomness.
(In particular, the bits they output will necessarily be ``very far''
from truly random, in a statistical sense.)  One easy-to-imagine
application would be to use the pseudorandom bit stream as a one-time
encryption pad, which would allow the shared secret key to be much
smaller than the message.  But what does it \emph{mean} for a string
of bits to be ``random-looking''?  And how can we be confident that
using such bits does not introduce any unforeseen weaknesses in our
system?

More generally, the motivating question for our study is:
\begin{center}
  When can two (possibly different) objects be considered
  \emph{effectively the same}?
\end{center}
The answer:
\begin{center}
  \emph{When they can't be told apart!}
\end{center}
Though seemingly glib, this answer encapsulates a very powerful
mindset that will serve us well as we go forward.

\subsection{Statistical Indistinguishability}
\label{sec:stat-indist}

We use probability theory to model (in)distinguishability.  If two
distributions are identical, then they certainly should be considered
indistinguishable.  We relax this condition to define
\emph{statistical} indistinguishability, for when the
\emph{statistical distance} between the two distributions is
negligible.  The statistical distance between two distributions $X$
and $Y$ over a domain $\Omega$ is defined as\footnote{For a set~$S$ of
  real numbers, its \emph{supremum} $\sup(S)$ is the least upper bound
  of the elements of~$S$. The suprenum can be thought of as a
  generalization of the maximum element to (infinite) sets where no
  such element may exist. For example, there is no maximum element of
  the set $[0,1) \subseteq \R$, but the suprenum is 1 since it is the
  smallest real number that is larger than every element of
  $[0,1)$. When~$S$ is finite, the supremum is simply the maximum.}
\[ \Delta(X,Y) := \sup_{A \subseteq \Omega} \abs{X(A) - Y(A)}, \]
where $X(A) = \sum_{w \in A} \Pr[X=w]$ is the probability that a draw
from~$X$ lands in~$A$, and likewise for~$Y(A)$. We can
view~$A \subseteq \Omega$ as a statistical ``test'' that has some
probability of ``passing'' when given an element drawn from~$X$, or
from~$Y$; the statistical distance is essentially the maximum
difference between these two probabilities, taken over all tests.
Note that $A$ and $\bar{A}$ are effectively the same test, since
\[ \abs{X(\bar{A})-Y(\bar{A})} = \abs{1-X(A)-(1-Y(A))} =
  \abs{Y(A)-X(A)} = \abs{X(A)-Y(A)}. \]

\begin{lemma}
  \label{lem:stat-dist}
  For distributions $X,Y$ over a finite domain $\Omega$,
  \[ \Delta(X,Y) = \frac12 \sum_{w \in \Omega} \abs{X(w) - Y(w)}. \]
\end{lemma}

\begin{proof}
  Let the test $A = \set{w \in \Omega\; :\; X(w) > Y(w)}$.  This makes
  $X(A) - Y(A)$ as large as possible, so
  $\Delta(X,Y) = X(A)-Y(A) = \sum_{w \in A} \abs{X(w)-Y(w)}$.  As
  noted above, we also have
  $\Delta(X,Y) = Y(\bar{A})-X(\bar{A}) = \sum_{w \in \bar{A}}
  \abs{X(w)-Y(w)}$.  Summing the two equations, we have
  $2\Delta(X,Y) = \sum_{w \in \Omega} \abs{X(w)-Y(w)}$, as desired.
\end{proof}

Statistical distance is very robust, which enhances its usefulness.
Using \cref{lem:stat-dist}, the following facts are
straightforward to prove.

\begin{lemma}
  \label{lem:ipl}
  Let $f$ be a function (or more generally, randomized procedure) on
  the domain of $X,Y$. Then $\Delta(f(X), f(Y)) \leq \Delta(X,Y)$.
\end{lemma}
In other words, statistical distance cannot be increased by the
application of any (randomized) procedure.

\begin{lemma}
  Statistical distance is a metric; in particular, it satisfies the
  following three properties:
  \begin{itemize}[itemsep=0pt]
  \item Identity of indiscernibles: $\Delta(X,Y) = 0 \iff X = Y$
  \item Symmetry: $\Delta(X,Y) = \Delta(Y,X)$
  \item Subadditivity (triangle inequality):
    $\Delta(X,Z) \leq \Delta(X,Y) + \Delta(Y, Z)$
  \end{itemize}
\end{lemma}

\begin{question}
  Justify why it is the case that $\Delta(X,Y) = 0 \iff X = Y$.
\end{question}

\begin{answer}
  In the forward direction, because $\Delta(X,Y) = 0$ and all the
  terms in the sum from \cref{lem:stat-dist} are non-negative,
  we have that $\abs{X(w) - Y(w)} = 0$ for all $w \in \Omega$, i.e.,
  $X(w) = Y(w)$. Hence, $X = Y$. In the other direction, if $X = Y$,
  then by definition $X(A) - Y(A) = 0$ for all $A$, so
  $\Delta(X,Y) = \sup(\set{0}) = 0$.
\end{answer}

Statistical distance lets us say when two (sequences of) distributions
are ``essentially the same,'' in an asymptotic sense.

\begin{definition}
  \label{def:stat-ind}
  Let $\calX = \set{X_{n}}$ and $\calY = \set{Y_{n}}$ be sequences of probability distributions, called \emph{ensembles}, indexed by security parameters $n\in\N$.\footnote{More generally, to accommodate restrictions on the allowed security parameters, both ensembles may be indexed by a common infinite subset of~$\N$.
    In that case, all asymptotic statements refer to~$n$ tending to infinity within that subset.}
  We say that $\calX$ and $\calY$ are \emph{statistically indistinguishable}, written $\calX \statind \calY$, if \[ \Delta(X_{n}, Y_{n}) = \negl(n)\text. \]
\end{definition}

\begin{example}
  Let $X_{n}$ be the uniform distribution over $\bit^{n}$, and let
  $Y_{n}$ be the uniform distribution over the nonzero strings
  $\bit^{n} \backslash \set{0^{n}}$.  An optimal test $A$ is the
  singleton set $A = \set{0^{n}}$, yielding
  $\Delta(X_{n}, Y_{n}) = 2^{-n} = \negl(n)$, so
  $\calX \statind \calY$.  (This can also be seen by calculating the
  summation in \cref{lem:stat-dist}.)  The analysis extends
  similarly to any $Y_{n}$ that leaves out a $\negl(n)$ fraction of
  $\bit^{n}$.  From this we can say that such ensembles $\calY$ are
  ``essentially uniform,'' or \emph{statistically pseudorandom}.
\end{example}

Can a statistically pseudorandom generator exist?  This
depends on the definition of ``generator'' (which we give in the next lecture), but
for any meaningful definition of the term, it isn't possible!  This
can be shown by explicitly demonstrating an appropriate subset (or
test) that distinguishes strings output by the generator from
uniformly random ones; we will do this in the next lecture.

\subsection{Computational Indistinguishability}
\label{sec:comp-indist}

We can define a natural analogue of statistical distance in the
computational setting, where the ``test'' is implemented by an
\emph{efficient algorithm}.  Namely, for distributions $X$ and $Y$ and
an algorithm $\Adv$ (possibly randomized), define $\Adv$'s
\emph{distinguishing advantage} between $X$ and $Y$ as
\[ \advan_{X,Y}(\Adv) = \abs*{\Pr[\Adv(X)=1] - \Pr[\Adv(Y)=1]}. \]
(The output of $\Adv$ can be arbitrary, but we interpret $1$ as a
special output indicating that the test implemented by $\Adv$ is
``satisfied,'' and any other output as ``not satisfied.'')
For ensembles $\calX = \set{X_{n}}$ and $\calY = \set{Y_{n}}$, we give the distinguisher the security parameter explicitly and define its advantage as the function
\[
  \advan_{\calX,\calY}(\Adv)(n)
  = \abs*{\Pr[\Adv(1^{n},X_{n})=1] - \Pr[\Adv(1^{n},Y_{n})=1]}\text.
\]
We may suppress the input $1^{n}$ and the argument~$n$ when they are clear from context.

\begin{definition}
  \label{def:comp-ind}
  Let $\calX = \set{X_{n}}$ and $\calY = \set{Y_{n}}$ be ensembles.
  We say that $\calX$ and $\calY$ are \emph{computationally indistinguishable}, written $\calX \compind \calY$, if $\advan_{\calX,\calY}(\Adv) = \negl(n)$ for every PPT algorithm~$\Adv$.

  \noindent If each~$X_{n}$ is a distribution over $\bit^{\ell(n)}$, we say that~$\calX$ is (computationally) \emph{pseudorandom} if $\calX \compind \set{U_{\ell(n)}}$, the ensemble of uniform distributions over $\bit^{\ell(n)}$.
\end{definition}

The basic facts about statistical distance also carry over to computational indistinguishability, where we restrict all functions/tests to be \emph{efficient}.
We call the analogue of \cref{lem:ipl} the \emph{composition lemma} (see \cref{lem:composition}), because it involves the composition of an arbitrary PPT procedure and distinguisher.
We call the analogue of the triangle inequality the \emph{hybrid lemma} (see \cref{lem:hybrid}), because we typically invoke it on a sequence of ``hybrid'' distributions that interpolate step-by-step between two distributions we ultimately care about.

\subsubsection{Composition Lemma}

\begin{lemma}[Composition lemma]
  \label{lem:composition}
  Let $\calX = \set{X_{n}}$ and $\calY = \set{Y_{n}}$ be ensembles, and let~$\Sim$ be a PPT algorithm.
  Define $X'_{n}=\Sim(1^{n},X_{n})$ and $Y'_{n}=\Sim(1^{n},Y_{n})$, and let $\calX'=\set{X'_{n}}$ and $\calY'=\set{Y'_{n}}$.
  For any distinguisher~$\Dist'$ (regardless of efficiency), the composed distinguisher
  \[
    \Dist(1^{n},\cdot) := \Dist'(1^{n},\Sim(1^{n},\cdot))
  \]
  has advantage
  \[
    \advan_{\calX,\calY}(\Dist)(n) = \advan_{\calX',\calY'}(\Dist')(n)\text.
  \]
  If~$\Dist'$ is PPT, then so is~$\Dist$.
  In particular, if $\calX \compind \calY$, then $\calX' \compind \calY'$.
\end{lemma}

\begin{proof}
  The algorithm~$\Dist$ runs~$\Sim$ and then~$\Dist'$, using independent randomness for the two runs, and outputs the answer of~$\Dist'$.
  If~$\Dist'$ is PPT, then so is~$\Dist$, because~$\Sim$ is PPT and its output length is bounded by its running time.
  By construction,
  \[
    \begin{aligned}
      \advan_{\calX,\calY}(\Dist)(n)
      &= \abs*{\Pr[\Dist(1^{n},X_{n})=1] - \Pr[\Dist(1^{n},Y_{n})=1]} \\
      &= \abs*{\Pr[\Dist'(1^{n},\Sim(1^{n},X_{n}))=1]
        - \Pr[\Dist'(1^{n},\Sim(1^{n},Y_{n}))=1]} \\
      &= \abs*{\Pr[\Dist'(1^{n},X'_{n})=1] - \Pr[\Dist'(1^{n},Y'_{n})=1]} \\
      &= \advan_{\calX',\calY'}(\Dist')(n)\text.
    \end{aligned}
  \]
  For the final claim, restrict to PPT distinguishers~$\Dist'$.
  If $\calX \compind \calY$, the left-hand side is negligible because~$\Dist$ is PPT, so the right-hand side is negligible as well; since~$\Dist'$ is an arbitrary PPT algorithm, $\calX' \compind \calY'$.
\end{proof}

The algorithm~$\Sim$ is often called a \emph{simulator} because it ``simulates'' a sample from~$X'_n$ or~$Y'_n$ given a sample from~$X_n$ or~$Y_n$, respectively.
It is thus a \emph{reduction} from distinguishing $\calX$ and~$\calY$ to distinguishing $\calX'$ and~$\calY'$.

\begin{question}
  Let $\calX=\set{X_n}$ be a pseudorandom ensemble.
  Is $\set{\overline{X_n}}$ necessarily pseudorandom as well, where $\overline{\cdot}$ denotes bitwise complement?
  \emph{Hint:} use \cref{lem:composition}.
\end{question}

\begin{answer}
  Yes: complementing a uniformly random string preserves its distribution.
  Let $\ell(n)$ denote the length of the strings distributed according to~$X_n$.
  Define the efficient algorithm $\Sim(1^n,z):=\overline{z}$.
  Since $\calX\compind\set{U_{\ell(n)}}$, \cref{lem:composition} gives
  \[
    \set{\overline{X_n}}
    =\set{\Sim(1^n,X_n)}
    \compind\set{\Sim(1^n,U_{\ell(n)})}
    =\set{U_{\ell(n)}}\text,
  \]
  where the last equality of distributions holds because bitwise complementation is a permutation of $\bit^{\ell(n)}$.
  Thus $\set{\overline{X_n}}$ is pseudorandom.
\end{answer}

\subsubsection{Hybrid Lemma}

\begin{lemma}[Hybrid lemma]
  \label{lem:hybrid}
  Let $X^0,\ldots,X^m$ be distributions for some positive integer~$m$, and let~$\Dist$ be any algorithm (regardless of efficiency).
  Writing $p_i=\Pr[\Dist(X^i)=1]$, we have
  \[
    \advan_{X^0,X^m}(\Dist)
    =\abs[\bigg]{\sum_{i\in[m]}\parens{p_{i-1}-p_i}}
    \leq\sum_{i\in[m]}\advan_{X^{i-1},X^i}(\Dist)\text.
  \]
\end{lemma}

\begin{proof}
  By telescoping and the triangle inequality,
  \[
    \advan_{X^0,X^m}(\Dist)=\abs{p_0-p_m}=\abs[\bigg]{\sum_{i\in[m]}\parens{p_{i-1}-p_i}}\leq\sum_{i\in[m]}\abs{p_{i-1}-p_i}=\sum_{i\in[m]}\advan_{X^{i-1},X^i}(\Dist)\text.
  \]
\end{proof}

It is important to understand that even though the intended goal of~$\Dist$ is to distinguish~$X^0$ from~$X^m$, the proof of \cref{lem:hybrid} analyzes the same distinguisher on \emph{every} hybrid distribution~$X^i$.
This is legitimate because~$\Dist$ is just an algorithm, so we can analyze its behavior on any input distribution we like.
If~$\Dist$ has advantage~$\varepsilon$ in distinguishing between~$X^0$ and~$X^m$, the hybrid lemma guarantees that it has advantage at least $\varepsilon/m$ in distinguishing between some adjacent pair $X^{i-1}$ and~$X^i$.

\begin{corollary}
  \label{cor:hybrid-constant}
  Let $\calX^0,\calX^{1},\ldots,\calX^m$ be ensembles, where~$m$ is a positive constant independent of~$n$.
  If $\calX^{i-1}\compind\calX^i$ for every $i\in[m]$, then $\calX^0\compind\calX^m$.
\end{corollary}

\begin{proof}
  Write $\calX^i=\set{X^i_n}$ and fix any PPT distinguisher~$\Dist$.
  For each security parameter~$n$, apply \cref{lem:hybrid} to $X^0_n,\ldots,X^m_n$ and the algorithm $\Dist(1^n,\cdot)$, obtaining
  \[
    \advan_{\calX^0,\calX^m}(\Dist)(n)
    \leq\sum_{i\in[m]}\advan_{\calX^{i-1},\calX^i}(\Dist)(n)\text.
  \]
  Each summand is a negligible function by hypothesis, and a sum of a constant number of negligible functions is itself negligible: for every constant $c>0$, each summand is $O(n^{-c})$, so their sum is also $O(n^{-c})$.
  Since~$\Dist$ was arbitrary, $\calX^0\compind\calX^m$.
\end{proof}

\paragraph{Polynomially many hybrid steps.}

In many applications, the number of hybrid steps grows polynomially with the security parameter~$n$.
We cannot directly express this by replacing the constant~$m$ in \cref{cor:hybrid-constant} with a function $m(n)$, because~$n$ is in scope only within each ensemble.
Instead, for each~$n$, we define a sequence of distributions $H^0_n,\ldots,H^{m(n)}_n$; the initial and final distributions form the two ensembles $\calH^0:=\set{H^0_n}$ and $\calH^m:=\set{H^{m(n)}_n}$ that we want to prove are indistinguishable.
The hybrid bound (\cref{lem:hybrid}) applies to each such sequence, but bounding this sum by a negligible function requires additional care: the sum of polynomially many \emph{different} negligible functions is not necessarily negligible.\footnote{For example, let $\nu_i(n)=1$ when $n=i$ and $\nu_i(n)=0$ otherwise.
  Each~$\nu_i$ is negligible, but $\sum_{i=1}^n\nu_i(n)=1$ for every~$n$.}

\begin{corollary}[Polynomially many hybrid steps]
  \label{cor:hybrid-polynomial}
  Let $m(n)\geq1$ be integer-valued, with $1^{m(n)}$ polynomial-time computable from $1^n$, so $m(n)=\poly(n)$.
  For each~$n$, let $H^0_n,\ldots,H^{m(n)}_n$ be distributions, and define $\calH^0:=\set{H^0_n}$ and $\calH^m:=\set{H^{m(n)}_n}$.
  Let $\calX=\set{X_n}$ and $\calY=\set{Y_n}$ be computationally indistinguishable ensembles, and let~$\Sim$ be a PPT algorithm.
  Writing $\Sim_i(1^n,z):=\Sim(1^n,i,z)$, suppose that for every~$n$ and every $i\in[m(n)]$,
  \begin{equation}
    \label{eq:hybrid-simulator-identities}
    \Sim_i(1^n,X_n)\equiv H^{i-1}_n
    \qquad\text{and}\qquad
    \Sim_i(1^n,Y_n)\equiv H^i_n\text.
  \end{equation}
  Then $\calH^0\compind\calH^m$.
\end{corollary}

\begin{proof}
  Define the \emph{averaging simulator}~$\bar{\Sim}$ that, on input $(1^n,z)$, chooses $i\gets[m(n)]$ uniformly and runs $\Sim_i(1^n,z)$ with fresh randomness; clearly, this~$\bar{\Sim}$ is PPT\@.

  Fix any PPT algorithm~$\Dist'$; we will show that it has negligible advantage in distinguishing~$\calH^0$ and~$\calH^m$.
  Let \[ p_i(n)=\Pr[\Dist'(1^n,H^i_n)=1] \text, \] so the advantage is $\abs{p_0(n)-p_{m(n)}(n)}$.
  To use the hypothesis $\calX\compind\calY$, consider the composed PPT distinguisher $\Dist(1^n,\cdot):=\Dist'(1^n,\bar{\Sim}(1^n,\cdot))$ from \cref{lem:composition}.
  Unfolding the definition of~$\Dist$, averaging over the choice of~$i$ using \cref{eq:hybrid-simulator-identities}, and telescoping the sum as in \cref{lem:hybrid} gives
  \begin{align*}
    \advan_{\calX,\calY}(\Dist)(n)
    &=\abs*{\Pr[\Dist'(1^n,\bar{\Sim}(1^n,X_n))=1]-\Pr[\Dist'(1^n,\bar{\Sim}(1^n,Y_n))=1]} \\
    &=\frac{1}{m(n)}\abs[\bigg]{\sum_{i=1}^{m(n)}\parens*{\Pr[\Dist'(1^n,H^{i-1}_n)=1]-\Pr[\Dist'(1^n,H^i_n)=1]}} \\
    &=\frac{1}{m(n)}\abs[\bigg]{\sum_{i=1}^{m(n)}\parens{p_{i-1}(n)-p_i(n)}} \\
    &=\frac{\abs{p_0(n)-p_{m(n)}(n)}}{m(n)}
    =\frac{\advan_{\calH^0,\calH^m}(\Dist')(n)}{m(n)}\text.
  \end{align*}
  Since $\calX\compind\calY$ and~$\Dist$ is PPT, its advantage on the left-hand side is negligible.
  Multiplying by $m(n)=\poly(n)$ preserves negligibility, so the advantage of~$\Dist'$ on the right-hand-side is negligible as well.
  Since~$\Dist'$ was arbitrary, we finally conclude that $\calH^0\compind\calH^m$, as claimed.
\end{proof}

What matters in \cref{eq:hybrid-simulator-identities} is that the mapping of $X_n,Y_n$ to the two adjacent hybrids has the same orientation for every~$i$.
That is, mapping~$X_n$ to $H^i_n$ and~$Y_n$ to $H^{i-1}_n$ for every~$i$ would work equally well: it simply negates the telescoping sum, leaving its absolute value unchanged.
Either consistent orientation makes the acceptance-probability differences telescope \emph{before} we take the absolute value.
By contrast, averaging the absolute values of the differences $p_{i-1}(n)-p_i(n)$ would only \emph{upper}-bound the advantage of~$\Dist$ by the average advantage of~$\Dist'$ between adjacent hybrids.
Such a bound goes in the wrong direction: we need to upper-bound the advantage of~$\Dist'$ between~$\calH^0$ and~$\calH^m$ in terms of the advantage of~$\Dist$ between~$\calX$ and~$\calY$, which we know is negligible.

\subsubsection{Changing the Security Parameter}

A cryptographic primitive evaluated at one security parameter may invoke another primitive at a different parameter.
For example, in Lecture~3's hardness amplification, the direct product~$f^{(m)}$ has input length $n'=n\cdot m(n)$, but evaluates~$f$ on blocks of length~$n$.
There we transferred a security guarantee expressed in~$n$ to one expressed in~$n'$.
Here we ask the analogous question for computational indistinguishability: if two ensembles are indistinguishable at parameter~$n$, do they remain so when indexed by~$n'$?

\begin{lemma}[Change of security parameter]
  \label{lem:change-parameters}
  Let $I\subseteq\N$ be infinite, and let $\calX = \set{X_{n}}_{n\in I}$ and $\calY = \set{Y_{n}}_{n\in I}$ be ensembles.
  Let $n'=n'(n)$ be an injective map from~$I$ to~$\N$ such that $1^{n'}$ is computable in polynomial time from $1^{n}$ for $n\in I$, and let $I'=n'(I)$.
  Define $X'_{n'(n)}=X_n$ and $Y'_{n'(n)}=Y_n$, and let $\calX'=\set{X'_{n'}}_{n' \in I'}$ and $\calY'=\set{Y'_{n'}}_{n' \in I'}$.
  If $\calX \compind \calY$, then $\calX' \compind \calY'$ with respect to the security parameter $n'\in I'$.
\end{lemma}

\begin{proof}
  Injectivity ensures that $\calX',\calY'$ are well-defined.
  Fix any PPT distinguisher~$\Dist'$ for the relabeled ensembles, and define
  \[
    \Dist(1^n,z):=\Dist'(1^{n'(n)},z)\text.
  \]
  Then~$\Dist$ is PPT because $1^{n'(n)}$ is efficiently computable from $1^n$.
  By construction, for every $n\in I$,
  \[
    \advan_{\calX,\calY}(\Dist)(n)
    =\advan_{\calX',\calY'}(\Dist')(n'(n))\text.
  \]
  Because $\calX\compind\calY$, this common advantage is negligible in~$n$ for $n\in I$.
  Efficient computation of $1^{n'}$ from $1^n$ implies $n'=\poly(n)$, and injectivity implies $n'(n)\to\infty$ as $n\to\infty$ within~$I$.
  The change-of-parameter property from Lecture~2 therefore makes the advantage negligible in~$n'$ for $n'\in I'$.
  Since~$\Dist'$ was arbitrary, $\calX'\compind\calY'$.
\end{proof}

\begin{remark}[Non-uniform adversaries]
  The composition lemma (\cref{lem:composition}), hybrid lemma and its corollaries (\cref{lem:hybrid,cor:hybrid-constant,cor:hybrid-polynomial}), and change-of-parameter lemma (\cref{lem:change-parameters}) have non-uniform analogues.
  In each, the indistinguishability hypotheses and conclusions alike are against non-uniform PPT distinguishers, and the simulators may also be non-uniform.
  For composition, give the reduction both the simulator's and the distinguisher's advice; their combined length is still polynomial.
  In the hybrid arguments, use the same advice on every hybrid distribution, so the same probability calculations apply.
  For a change of parameter, the distinguisher's advice at~$n'$ has polynomial length in~$n$ because $n'=\poly(n)$.
\end{remark}

\end{document}

%%% Local Variables: 
%%% mode: latex
%%% TeX-master: t
%%% End: 
